\chapter{Introduction}\label{chap:01-introduction}

\section{Two curves, one program}\label{sec:intro-conferences}

\paragraph{Vancouver, spring 2018}
Joelle Pineau puts two learning curves on screen. A simulated cheetah learns to run. Each curve
follows its score over training, averaged over five runs. The two curves separate clearly and a test
declares them different: $t = -9.09$, $p = 0.0016$. A reviewer would see two algorithms. Yet both
curves come from the same program, the TRPO algorithm, with the same settings. Only the random seed
changes, the starting number from which the computer draws all the randomness of a training run
\parencite{henderson2018matters}. The scene takes place during the invited talk at ICLR
(\emph{International Conference on Learning Representations}) in Vancouver \parencite{pineau2018iclr},
titled \enquote{Reproducibility, Reusability, and Robustness in Deep Reinforcement Learning}. Pineau
is then a professor at McGill University in Montreal and head of Facebook's artificial intelligence
research lab in that city. In December, she takes up and extends her argument at NeurIPS
(\emph{Neural Information Processing Systems}) in Montreal, as the Posner Lecture:
\enquote{Reproducible, Reusable, and Robust Reinforcement Learning} \parencite{pineau2018neurips}.
ICLR, NeurIPS and ICML (\emph{International Conference on Machine Learning}) are the three main
conferences of machine learning.

\paragraph{The starting question}
When one team publishes a result, does another team that follows the recipe get the same one? Pineau
starts from a survey by the journal \emph{Nature}: of 1,576 researchers across all disciplines,
more than 70\% said they had already failed to reproduce a colleague's experiment, 52\% saw a
\enquote{significant} reproducibility crisis and 38\% a \enquote{slight} one
\parencite{baker2016reproducibility}. At a workshop in December 2017, she had put the same question to
twenty machine learning specialists, 80\% of whom saw a crisis, slight or significant
\parencite{pineau2018iclr}. She then quotes \textcite{buckheit1995wavelab}, for whom a scientific
article about computation is only the advertising of the research. The research itself is the complete
software environment and the full set of instructions that produced the figures.

\paragraph{A growing field, costly errors}
Pineau recalls that for nearly twenty years RL cited a single landmark success, the backgammon program
TD-Gammon \parencite{tesauro1995tdgammon}, before AlphaGo \parencite{silver2016alphago}. Around 2000,
an RL workshop gathered thirty-five people. For 2016, she counts about thirteen thousand RL papers
published in the year \parencite{pineau2018iclr}. She adds that RL is approaching problems that
matter: the design of medical treatments, the management of natural resources, of energy networks and
of financial systems. Practitioners doubted too. In February 2018, Alex Irpan, of Google
Brain, wrote that he answered no to anyone who asked him whether RL could solve their problem. He
believed he was right at least 70\% of the time \parencite{irpan2018rlhard}. Pineau puts the address of
this blog post on screen at NeurIPS.

\paragraph{The sixteen flaws}
Pineau then shows, figure after figure, what RL lacks to deserve trust. Most of the figures come from
her team's paper signed by Peter Henderson and co-authors \parencite{henderson2018matters}. We recorded
sixteen flaws in the two talks, labelled F1 to F16 and sorted into four families: the materials needed
to rerun an experiment (code, plan, software), the tuning of hyperparameters (the settings an
experimenter fixes by hand) and of baselines (the existing methods a new one is compared with), random
seeds (the integer from which a training run draws all its randomness) and statistics, the reach of the
conclusions. The most famous is flaw F11, the two curves of the same program.
Chapter~\ref{chap:02-contexte} describes the sixteen flaws one by one
(Table~\ref{tab:failles}). Pineau also proposes two remedies: the ICLR 2018
reproducibility challenge, in which anyone could try to reproduce a submitted paper, then a
reproducibility checklist, which NeurIPS adopted in 2019 \parencite{pineau2021improving}.

\section{The three Rs: reproducible, reusable, robust}\label{sec:intro-trois-r}

\paragraph{What Pineau defines, what she shows}
The titles of the two talks name three requirements. Pineau defines only one of them, through a
quotation. The second page of her NeurIPS talk takes it from the report written by
\textcite{bollen2015robust} for the US National Science Foundation: reproducibility \enquote{refers to the
ability of a researcher to duplicate the results of a prior study\textellipsis{} using the same materials as were used
by the original investigator}. It is \enquote{a minimum necessary condition for a finding to be
believable and informative}. The full definition, on page~3 of that report, reads \enquote{the same
materials and procedures}. One slide in Vancouver ties the three words into one sentence (ICLR,
slide 16): \enquote{reusable material (software, datasets, experimental platforms) help us reproduce
scientific methods to establish the robustness of the findings}. The other two words are shown, not
defined. For reusability, Pineau shows the \enquote{reusable platforms for RL}, MuJoCo and the Arcade Learning
Environment among them (ICLR, slide 15), the call to \enquote{commit to sharing reusable material} (ICLR,
slide 44), then the lines of her checklist that ask for downloadable source code \enquote{including all
dependencies} and for the simulation environment (NeurIPS, pages 31 to 33). For robustness, she shows
the pages titled \enquote{Robustness of policy gradient algorithms}, where the ranking of four algorithms changes
from \task{HalfCheetah} to \task{Hopper} and \task{Swimmer} (NeurIPS, pages 14 and 15). Then come
implementations that diverge (pages 16 and 17), random seeds (pages 21 and 22), the move \enquote{from
fair comparisons\textellipsis{} to robust conclusions} (page 28) and benchmarks that are \enquote{brittle to
perturbations} (page 44).

\paragraph{Neighbouring vocabularies}
The report Pineau quotes separates reproduction, with the same data, from replication, which follows
the same procedures on new data. It calls a finding robust when it is reproducible, replicable and
generalisable to other contexts \parencite[p.~3]{bollen2015robust}. The report of the US National
Academies \parencite[Summary, p.~1]{nasem2019reproducibility} separates two words. Reproducibility there
means consistent results obtained with the same input data, computational steps, methods, code and
conditions of analysis. Replicability means consistent results across studies that answer the same
question, each with its own data. The report by \textcite[Section~2.1, Figure~1]{pineau2021improving}
crosses two axes, data and analysis, which give four cells:
\begin{itemize}
  \item same data and same analysis: the work is reproducible;
  \item new data and same analysis: it is replicable;
  \item same data and another analysis, such as rewritten code: it is robust;
  \item new data and another analysis: it is generalisable.
\end{itemize}
In RL, training produces its own data. A new random seed produces new
data. Rerunning the same code on other random seeds is therefore replication in all three vocabularies.
We keep Pineau's three words and state, for each, what is held fixed, what varies and what \enquote{the
same result} means.

\paragraph{The same result}
A learning curve changes from one machine to another: on another processor, MuJoCo, JAX and PyTorch
round differently, so the same random seed gives another curve (\cref{sec:outils}). The result that
reproduces is a conclusion drawn from several runs by a rule written in advance. \enquote{PPO JAX beats
PPO PyTorch on \task{Hopper}} is true or false under the settled-winner rule (\cref{sec:regle}),
applied to the scores of ten runs. Getting the same result means getting the same verdict, by the same
rule, from a new set of independent runs. The rule rests on a finite number of random seeds. It
therefore sometimes declares a winner between two groups of runs of the same program: \cref{tab:regle-tranche}
measures that rate. A single failed reproduction does not refute a conclusion, as
\textcite[p.~4]{bollen2015robust} already said of replications: two studies that seem to disagree
sometimes differ only by the sampling fluctuations of each. A conclusion drawn from three runs
reproduces poorly, as experiment E5 measures (\cref{sec:e5}).

\paragraph{Reproducible}
A result is reproducible when another person, with the same code, the same hyperparameters, the same
task in the same version, the same budget of steps and the same evaluation protocol, reaches the same
conclusion. Only the random seeds and the machine change. Reproduction is checked at two levels. The
first is the identical rerun: same random seed, same machine, a single computing thread. It gave back
the same evaluation points for PPO PyTorch and PPO JAX (\cref{sec:calcul}), and to the byte for three training
runs of SAC and CrossQ redone unchanged (\cref{sec:graines}). This level tests the code. The second level tests the conclusion: other
random seeds, the same verdict. Bollen's report, the National Academies report and the 2021 grid call this level a replication. Experiment
E1 measures how often two groups of random seeds of the same program pass for two algorithms
(\cref{sec:e1}). E5 counts the random seeds needed to recover the right ranking (\cref{sec:e5}). The
consolidation reruns E2, E3, E5, E7 and B3 on random seeds 10 to 19 (\cref{sec:consolidation}). The
survey measures whether a paper supplies what reproduction requires: the code (criterion Q5), the
hyperparameters of the method and of the baselines and how they were chosen (Q3, Q4, Q6), the number of
training runs (X1), the compute (X7), the versions of the environments and of the dependencies (Y7,
Y10), as listed in \cref{tab:criteres}.

\paragraph{Reusable}
A piece of work is reusable when others take its material (code, simulated environment, dataset,
reading grid) for another question without rewriting it, then get the behaviour it describes. The
material is held fixed. The user, the question, the date and the software around the material change.
Reusability qualifies material. It serves reproduction, in the order of slide 16. Experiment E6 shows
its limit (\cref{sec:e6}). The hyperparameters that the RL Baselines3 Zoo collection publishes for PPO
on \task{Hopper}, reused with the libraries of 2026, give a robot that stands upright without moving
forward. With the libraries of 2022, the same recipe gave a robot that runs. Reused without its
versions, this material changed its result. Our repository therefore freezes its environments with the hash of
every package (\cref{sec:envs}, \cref{chap:D-environnements}) and publishes the package and its file
formats (\cref{sec:paquet}, \cref{sec:formats}) under the BSD 3-Clause licence (\cref{sec:licence}).
The command \texttt{make verify} recomputes every analysis and compares the result files byte for byte
(\cref{sec:make}). The reading grid, with its decision rules, serves for other papers
(\cref{chap:A-grille}). In the survey, the link to the code (Q5) and the list of dependencies with
their versions (Y10) measure the reusable material of a paper.

\paragraph{Robust}
A conclusion is robust when it keeps its verdict under a change it claims not to depend on. The
algorithm (the published idea) and the statistical rule are held fixed. One factor changes at a time.
Our experiments test four factors, each in an experiment of its own: the task (E2, \cref{sec:e2}), the program that implements the same
algorithm (E3, PPO written in PyTorch and in JAX, \cref{sec:e3}), the performance measure (B3,
\cref{sec:b3}) and the physics of the evaluation world (E7 and the graded physics, \cref{sec:e7} and
\cref{sec:dose}). The survey asks whether papers run these tests: the number of tasks behind the
conclusion (Y3), a task where the method does not win (Y4), the implementation of the baselines (Y5),
an evaluation outside the training conditions (Q7, Y8), the sensitivity to one hyperparameter (Y9).
This sense of \enquote{robust} covers what the 2021 grid calls robust (E3 changes the code) and what it
calls generalisable (E2 and E7 change the task or the physics). It matches the sense of
\textcite{bollen2015robust}, for whom a robust finding reproduces, replicates and generalises. It also
follows Pineau's pages, which put the change of task under the title of robustness.

\paragraph{The order of the three requirements}
The three requirements follow one another in the order of slide 16. Reusable material makes
reproduction possible. A conclusion reproduced on other random seeds then becomes the reference against
which the same conclusion is compared under another task, another program or another physics.

\section{The question}\label{sec:intro-question}

Eight years have passed. RL has changed in size and in terrain: in 2025 and 2026, 27 of the 50 RL papers
in the first random draw of our survey apply what they call RL to large language models. This
report asks one question: \emph{did we listen to Joelle Pineau?} For each of the sixteen flaws, we ask
whether the papers published between 2018 and 2026 still show it, then what it costs when the
experiments of 2018 are rerun with the algorithms of 2026. Pineau's list becomes a grid, applied with
the same severity to published papers and to our own experiments.

\section{Two instruments: reread the papers, rerun the experiments}\label{sec:intro-instruments}

Both instruments follow three protocols written before the work and published in the repository with
every change made to them afterwards (\file{study/PROTOCOL.md}, \file{sim/SIM\_PROTOCOL.md},
\file{sim/llm/LLM\_PROTOCOL.md}).

\paragraph{The survey}
The first instrument rereads papers with Pineau's reading grid (criteria Q1 to Q8), extended by seven
criteria drawn from her ICLR talk (X1 to X7) and by a second grid of twelve criteria for the flaws she
left unmeasured (Y1 to Y12). We had AI agents read a first draw of 50 RL papers from 2025 and 2026, out
of the 1,252 RL papers of ICLR 2026, ICML 2026 and NeurIPS 2025, then 50 classical RL papers from 2026
(robotics, games, control, language models excluded), 50 papers from 2018 and 20 papers per year for
2019, 2020, 2022 and 2024. These agents are programs built on a generative AI model that read each PDF
and fill in the grid (\cref{sec:vocabulaire}). Each \enquote{yes} rests on an exact
quotation and its page. A \enquote{no} says what was read without finding it. The verdict compares classical RL in 2026 with our own reading of 2018, using
the same grids and agents of the same model. The readers of 2018 and those of 2026 were different agents
(Section~\ref{sec:sm-annees}). Eight fresh agents therefore reread the experimental papers behind the
verdicts, each of them reading both years, to separate the effect of the year from that of the readers: three of the four verdicts hold, while the fourth keeps its direction without being established (Section~\ref{sec:sr-annees}).
Chapters~\ref{chap:03-sondage-methode}
and~\ref{chap:04-sondage-resultats} describe the method and the results.

\paragraph{Why agents read}
AI agents read the papers. Facts that anyone can check in the repository support that choice. Each
agent reads the whole paper, appendices included. The NeurIPS appendices shipped as archives were
extracted, then reread by each reader. Each \enquote{yes} quotes the exact sentence of the paper and
its page. A script finds every quote in the PDF: of the 89 quotes of the first grid not found at
first, none was invented. A \enquote{no} quotes nothing. The 377 final \enquote{no} decisions on the
verdict criteria were therefore set against the text with patterns fixed in advance. For 283 of them,
an arbiter judged the passage found. Ten turned into \enquote{yes} and no verdict changed. The
instructions given to each agent are published word for word (\cref{sec:sm-lecteurs}). The double
readings are made by copies of the same model. A reader from another model family,
\texttt{gemma3:12b}, therefore reread the 98 papers behind the verdicts. It finds the direction of all
four rises. Its disagreement with our readers does not depend on the year in any detectable way
(\cref{sec:sr-gemma}).

\paragraph{The experiments}
The second instrument reruns the experiments of 2018: four programs (PPO in PyTorch and in JAX, SAC, CrossQ) on four MuJoCo tasks, with ten
random seeds per program and per task, then twenty. Thirty PPO runs and thirty TRPO runs on
\task{HalfCheetah} complete the set, with one evaluation on random seeds never seen in training and
another under modified physics. A further strand illustrates the same flaws on language models trained
with GRPO. Chapters~\ref{chap:05-experiences-methode}, \ref{chap:06-experiences-resultats}
and~\ref{chap:07-modeles-de-langue} describe these experiments.

\section{What this report establishes}\label{sec:intro-contributions}

\paragraph{The findings}
\begin{enumerate}
  \item \textbf{Four practices changed in classical RL, all around the result.} On the same grid,
  linking the code rose from 31\% to 74\% of papers, showing the effect of a setting from 44\% to
  80\%, quantifying the compute from 27\% to 58\% and reporting an uncertainty from 63\% to 88\%. The
  first three rises survive our checks, among them a rereading of the 98 papers by eight fresh readers
  who each read both years. The fourth keeps its direction without being established: its verdict
  depends on two papers read partly by another model (\cref{sec:sr-principale,sec:sr-annees}).
  \item \textbf{The way the result is judged has barely moved.} The statistical test goes from 6\% to
  12\%, baselines described with all their settings from 35\% to 29\%, held-out evaluation from 33\%
  to 45\%, benchmark versions from 23\% to 26\%. None of these gaps is detectable: with fifty papers
  per year, the survey could see with 80\% power only a rise of 18 to 30 points depending on the
  criterion (\cref{sec:sr-principale,sec:val-stat}).
  \item \textbf{The error bar spread faster than the bar computed over random seeds.} Papers that
  compute their uncertainty over independent training runs go from 35\% to 42\% only, with no
  detectable change (criterion X8, \cref{sec:sr-principale}).
  \item \textbf{The rise is slow, with no jump after the 2019 NeurIPS checklist.} Twenty papers per
  year in 2019, 2020, 2022 and 2024 trace a steady climb (\cref{sec:sr-trajectoire}).
  \item \textbf{Chance fools the eye, sorting fools the test.} Among thirty PPO runs, two groups of
  five drawn at random fool Welch's test in 1.2\% of draws, the five best against the five worst in
  31\%. Keeping the best three runs out of ten inflates the mean by 33\% for this PPO, by 32\% at the
  median over the sixteen (program, task) pairs (\cref{sec:e1,sec:e4}).
  \item \textbf{The $t$ of 9 from 2018 has a probable explanation.} The script Peter Henderson
  published in 2017 measures the spread between episodes and ignores the spread between random seeds.
  Applied to our thirty TRPO runs, it exceeds 9.09 in 38\% of random draws. The $(t, p)$ pair of 2018
  rules out the other ways of computing more than it designates this script (\cref{sec:b1}).
  \item \textbf{The best program depends on the task, the budget and the physics.} At equal budget
  with twenty random seeds, CrossQ leads on all four tasks. Under physics modified by 20\%, PPO JAX
  drops from first to last place on \task{Hopper} (\cref{sec:e2,sec:e7}).
  \item \textbf{A published recipe ages with its libraries.} The RL Baselines3 Zoo settings for PPO on
  \task{Hopper}, reused in 2026, give robots that stand upright without moving forward (0.00~m/s).
  The agent the Zoo published in 2022 with these settings reaches a score of \num{2439} under our
  protocol. With a discount factor of 0.99 instead of 0.999, the recipe makes the robot run again, at
  \num{2649} (\cref{sec:e6}).
  \item \textbf{On small language models, no GRPO gain is established at five random seeds.} Qwen's
  apparent gain comes from shorter answers: the share of cut-off answers falls from 75\% to 6\%, while
  the accuracy of finished answers does not rise (\cref{sec:llm-resultats}).
\end{enumerate}

\paragraph{What is new in the method}
\begin{enumerate}
  \item Pineau's grid applied again on equal terms: both years read with the same rules by agents
  of the same model, reread by eight readers who each read both years, then by a reader from another
  model family (\cref{chap:03-sondage-methode}).
  \item A second grid of twelve criteria (Y1 to Y12), with its decision rules, for the
  flaws the 2018 grid did not measure (\cref{chap:A-grille}).
  \item Checks published with their results: thirty papers reread blind with an interval on each
  agreement, the quotations audited, each negative decision set against the text and, for each
  criterion, the smallest difference the survey could detect (\cref{chap:04-sondage-resultats}).
  \item A repository that checks itself: \texttt{make verify} recomputes every analysis from the
  result files, then checks compare every number of the summary document
  (\file{docs/conclusions-en.md}) and of the blog article with those files (\cref{chap:08-logiciel}).
\end{enumerate}

\section{Reading guide}\label{sec:intro-guide}

A reader in a hurry can read the abstract, Section~\ref{sec:intro-contributions},
Table~\ref{tab:failles}, Section~\ref{sec:sr-principale}, Figures~\ref{fig:grille-egale}
and~\ref{fig:e1-faux-positifs}, then Chapter~\ref{chap:11-conclusion}.

Chapter~\ref{chap:02-contexte} sets out the formal framework of RL, describes the programs and the
tasks, inventories the sixteen flaws and places the work among related work.
Chapters~\ref{chap:03-sondage-methode} and~\ref{chap:04-sondage-resultats} give the method and the
results of the survey. Chapters~\ref{chap:05-experiences-methode} and~\ref{chap:06-experiences-resultats}
do the same for the MuJoCo experiments. Chapter~\ref{chap:07-modeles-de-langue} describes the GRPO
strand. Appendix~\ref{chap:08-logiciel} presents the software, from training to documents.
Chapter~\ref{chap:10-limites} gathers the limitations. Chapter~\ref{chap:11-conclusion} answers the
question, says what the measurements establish with a pointer to their chapters, then draws
recommendations.
The flaw-by-flaw verdict is in Section~\ref{sec:sr-failles}. The appendices give the complete rules of the grid (Appendix~\ref{chap:A-grille}),
additional tables (Appendix~\ref{chap:B-tableaux}), the list of papers read
(Appendix~\ref{chap:C-articles}) and the software environments (Appendix~\ref{chap:D-environnements}).

A reader who wants to redo a computation finds, in each caption,
the result file that holds it.
